\chapter{Networks, Models and This Exam}\label{ch:models}

\begin{outcomes}
By the end of this chapter you will be able to:
\begin{tight}
  \item \textbf{Describe} what a network does, in terms of the four things every
        network must solve: naming, addressing, forwarding and error handling.
  \item \textbf{Explain} the layered model well enough to use it as a
        troubleshooting tool, which is the only reason this exam still cares
        about it.
  \item \textbf{Identify} the devices in a small enterprise network and say what
        each one decides.
  \item \textbf{Navigate} the CCNA v2.0 blueprint and this book against it.
\end{tight}
\end{outcomes}

\begin{prereq}
Nothing. This is the first chapter, and the course assumes no prior networking.
\end{prereq}

\begin{keyterms}
host $\bullet$ node $\bullet$ link $\bullet$ frame $\bullet$ packet $\bullet$
segment $\bullet$ encapsulation $\bullet$ layer $\bullet$ protocol $\bullet$
switch $\bullet$ router $\bullet$ access layer $\bullet$ distribution layer
$\bullet$ broadcast domain $\bullet$ collision domain
\end{keyterms}

\section{What a network has to solve}

Two computers in the same room, joined by a cable, can exchange bytes. Everything
else in networking exists because that arrangement does not scale. Add a third
machine and you need a way to say \emph{which} machine you mean. Add a building
and you need a way to find a machine you cannot see. Add the internet and you
need a way to reach a machine whose owner you have never met, over equipment
owned by people you will never meet either.

\begin{conceptbox}[The four problems, and the four answers]
Every network technology in this book is an answer to one of four questions.
\begin{tight}
  \item \textbf{Who am I talking to?} $\rightarrow$ addressing. MAC addresses
        locally (\chref{ethernet}), IP addresses globally (\chref{ipv4}).
  \item \textbf{How do I find them?} $\rightarrow$ forwarding. Switches learn
        (\chref{ethernet}); routers decide (\chref{routingtable}).
  \item \textbf{How do we not talk over each other?} $\rightarrow$ media access
        and loop prevention (\chref{ethernet}, \chref{stp}).
  \item \textbf{What if it goes wrong?} $\rightarrow$ error detection,
        acknowledgement and retransmission (\chref{transport}).
\end{tight}
When you meet a new protocol, ask which of the four it is answering. It will be
one of them.
\end{conceptbox}

\begin{definitionbox}[Protocol]
An agreed set of rules for a conversation: what may be said, in what order, in
what format, and what to do when the expected reply does not arrive.

Two people who both know to say ``hello'', wait, then speak, have a protocol.
Two computers need theirs written down to the bit, because neither can improvise.
Almost every capitalised acronym in this book --- ARP, DHCP, OSPF, HSRP --- is one
such agreement, and each exists because somebody had one of the four problems
above.
\end{definitionbox}

\section{Layers, and why they survive}

A network stack is split into layers so that each layer can be changed without
rewriting the others. Wi-Fi replaced the cable without anyone rewriting the web.
That is the engineering argument, and it is real, but it is not why the layered
model matters to you this year.

\begin{conceptbox}[The model is a troubleshooting tool]
When something does not work, the layered model tells you \emph{what to test
next}. If a cable is unplugged, checking DNS is a waste of a minute you do not
have in the exam. Work up: is the link up? Does it have an address? Can it reach
its gateway? Can it resolve a name? Every troubleshooting question in this exam
--- and there are eight topics' worth --- is answered by walking the layers in
order.

The v2.0 blueprint no longer asks you to \emph{describe} the OSI model, as v1.1
did. It assumes you use it.
\end{conceptbox}

\begin{longtable}{@{}C{1.1cm}L{3.0cm}L{5.0cm}L{5.6cm}@{}}
\toprule
\rowcolor{primary!15}
\textbf{\#} & \textbf{Layer} & \textbf{What it decides} & \textbf{What breaks, and the symptom} \\
\midrule
\endhead
7--5 & Application & What the data means & Wrong URL, expired certificate; ``the
site is down but ping works'' \\
\rowcolor{lightbg}
4 & Transport & Which program on the host, and whether loss is repaired &
Blocked port; ``it pings but the application times out'' \\
3 & Network & Which network, and the path between networks & Wrong IP, wrong
mask, no route; ``I can reach my own subnet but nothing else'' \\
\rowcolor{lightbg}
2 & Data link & Which device on this link & Wrong VLAN, trunk down, spanning tree
blocking; ``the port is up but nothing passes'' \\
1 & Physical & Bits on the wire & Cable, transceiver, duplex or speed mismatch;
``the port is down'', or worse, up with errors climbing \\
\bottomrule
\end{longtable}

The bottom two rows are where this exam starts: topic \tp{1.1} is entirely
physical-layer diagnosis, and it is examined in \chref{physical}.

\begin{definitionbox}[Encapsulation, and the names of things]
As data moves down the stack each layer wraps it in a header. The wrapped unit
gets a different name at each layer, and using the right one is how you show you
know which layer you are discussing.
\begin{tight}
  \item Application data becomes a \term{segment} at layer 4 (TCP) or a
        \term{datagram} (UDP), when a port number is attached.
  \item That becomes a \term{packet} at layer 3, when IP addresses are attached.
  \item That becomes a \term{frame} at layer 2, when MAC addresses are attached.
  \item That becomes \term{bits} at layer 1.
\end{tight}
Switches forward frames. Routers forward packets. That single sentence explains
most of the difference between the two devices.
\end{definitionbox}

\section{The devices, and what each one decides}

\ccnafig{%
\begin{tikzpicture}[font=\scriptsize]
  \Host{P1}{(0,1.4)}{PC1}
  \Phone{PH}{(0,-1.4)}{Phone}
  \Sw{A1}{(3.2,0)}{Access}
  \Sw{A2}{(3.2,-2.8)}{Access}
  \Msw{D1}{(6.6,-1.4)}{Distribution}
  \Rtr{R1}{(10.0,-1.4)}{Router}
  \Fw{FW}{(13.0,-1.4)}{Firewall}
  \Cld{IN}{(15.8,-1.4)}{Internet}
  \Srv{S1}{(6.6,1.4)}{Servers}

  \plainwire{P1}{A1}
  \plainwire{PH}{A1}
  \wire[xconn]{A1}{D1}{trunk}
  \wire[xconn]{A2}{D1}{trunk}
  \plainwire{D1}{S1}
  \wire{D1}{R1}{routed}
  \plainwire{R1}{FW}
  \plainwire{FW}{IN}
\end{tikzpicture}}%
{A small enterprise network. Access switches connect end devices; a distribution
switch aggregates them and routes between VLANs; the router reaches other
networks; the firewall polices what leaves. Every chapter in Parts~II to~IV
configures one piece of this picture.}{smallent}

\begin{longtable}{@{}L{3.2cm}L{5.6cm}L{6.0cm}@{}}
\toprule
\rowcolor{primary!15}
\textbf{Device} & \textbf{The decision it makes} & \textbf{Where in this book} \\
\midrule
\endhead
\textbf{Switch} (layer 2) & ``Which port leads to this MAC address?'' Learns by
watching source addresses; floods when it does not know &
\chref{ethernet}, then all of Part~II \\
\rowcolor{lightbg}
\textbf{Router} & ``Which interface leads to this IP network?'' Never floods;
drops what it cannot route & Part~III \\
\textbf{Multilayer switch} & Both. Switches within a VLAN, routes between them
using an SVI & \chref{trunking}, \chref{routingtable} \\
\rowcolor{lightbg}
\textbf{Access point} & Bridges wireless clients onto a wired VLAN &
\chref{wireless} \\
\textbf{Firewall} & ``Is this flow permitted?'' Stateful, unlike an access list &
\chref{acl}, \chref{vpn} \\
\bottomrule
\end{longtable}

\begin{definitionbox}[Two domains worth keeping straight]
A \term{collision domain} is the set of ports that could transmit at the same
instant and interfere. Every switch port is its own collision domain, which is
why modern switched networks do not collide --- and why a \cmd{collisions}
counter that is climbing means something is badly wrong, usually a duplex
mismatch (\chref{physical}).

A \term{broadcast domain} is the set of devices that receive each other's
broadcasts. One VLAN is one broadcast domain. Splitting broadcast domains is the
entire purpose of \chref{vlans}, and joining them again is the purpose of
\chref{routingtable}.
\end{definitionbox}

\section{This exam, and this book}

\begin{examnote}[How the blueprint is built, and how to read it]
Cisco publishes twenty-nine numbered topics in five domains. The number is the
whole contract: if a technique is not in a topic statement, it is not directly
examinable, and if it is, the \emph{verb} tells you the depth expected.

\begin{tight}
  \item \emph{Describe} --- you must be able to explain it. Six topics.
  \item \emph{Interpret}, \emph{Validate}, \emph{Select}, \emph{Use},
        \emph{Manage} --- you must be able to read output or make a choice.
        Seven topics.
  \item \emph{Configure} --- you must produce the commands. Eight topics.
  \item \emph{Troubleshoot} or \emph{Diagnose} --- you must find and fix a fault.
        Eight topics.
\end{tight}

Sixteen of twenty-nine are configure-or-fix. That is the single most useful fact
about this exam, and it should change how you study: time at a command line
beats time with a highlighter, by a wide margin.
\end{examnote}

Every chapter of this book carries its topic numbers under the title. If you want
to revise by blueprint rather than by chapter, Appendix~\ref{app:coverage} maps
all twenty-nine topics to the chapters that cover them.

\begin{pitfall}
\begin{tight}
  \item \textbf{Studying a v1.1 book.} The exam number, 200-301, did not change
        when the blueprint did. A book that teaches QoS, NTP, REST APIs, JSON or
        Terraform in a CCNA context is written for the old blueprint. Check the
        version before you trust the contents.
  \item \textbf{Learning the OSI layers as a list to recite.} Nobody will ask you
        to name layer 6. They will show you a broken network and expect you to
        work upwards.
  \item \textbf{Assuming ``switch'' means layer 2.} Multilayer switches route.
        Most of the distribution layer in a real network is switches doing
        routing, and topic \tp{2.1} explicitly includes switch virtual
        interfaces.
\end{tight}
\end{pitfall}

\begin{breakfix}[Three reports, three different layers, and no commands yet.]
You have not configured anything and cannot yet read a \cmd{show} command. That is
deliberate: this first Break \& Fix is about \emph{placing} a fault, not fixing
one, because deciding which layer to investigate is the decision that all the
later chapters depend on.

Three reports arrive within an hour.

\begin{tight}
  \item[\textbf{A.}] ``The whole of the second floor is offline.'' Nobody there can
        reach anything, including each other.
  \item[\textbf{B.}] ``I can get to the file server but not the intranet site.''
        This user's colleague, at the next desk, has no trouble with either.
  \item[\textbf{C.}] ``Everything is slow this afternoon.'' It was fine this
        morning. Several people on different floors agree.
\end{tight}

\textbf{Place each one before reading on.} Which layer would you start at, and
what single question would you ask first?

\medskip
\textbf{A --- start at layers 1 and 2.} A whole floor failing together, including
traffic \emph{between} hosts on that floor, means the fault is common to all of
them and below the network layer --- an uplink, a switch, or power. Nothing about
this is an addressing or routing problem, because hosts on the same segment do not
need routing to reach each other. First question: does anything on that floor
reach anything else on that floor?

\textbf{B --- start at layer 7, then 3.} One user, one service, while a neighbour
on the same segment is fine. That rules out the cable, the switch and the VLAN
immediately --- they are shared, and they are working. Suspect name resolution or
the client's own configuration. First question: does the intranet site work for
them \emph{by IP address}? That one test splits ``name resolution'' from
``everything else'' (\chref{dns}).

\textbf{C --- do not start anywhere yet.} ``Slow'' is not a fault report, it is a
feeling, and it is the one report that most often wastes an afternoon. It could be
a duplex mismatch (layer 1), a loop (layer 2), a saturated uplink, or a slow server
that has nothing to do with the network. First question: slow doing \emph{what},
from \emph{where}, compared with \emph{when}? Get one measurable, repeatable case
before touching anything.

\medskip
\textbf{The point.} Reports A and B contain enough information to place the fault
within a layer or two before you log in to anything. Report C does not, and the
correct first action is to go and get that information rather than to start
guessing. \chref{tshootmethod} builds this into a method; for now, notice that the
useful question was different in all three cases, and that in none of them was it
``what was broken last time''.
\end{breakfix}

\begin{checkpoint}
\begin{tightnum}
  \item A user reports that a web page will not load. Ping to the server's IP
        address succeeds. Which layers have you just ruled out, and which is the
        first one you should now suspect?
  \item A switch port shows as up, but the \cmd{collisions} counter is
        increasing steadily. Which layer is at fault, and what is the most likely
        single cause?
  \item Why does a router drop a packet for an unknown network when a switch
        floods a frame for an unknown MAC address?
\end{tightnum}
\end{checkpoint}

\begin{chaptersummary}
\begin{tight}
  \item Networks solve four problems: addressing, forwarding, media access and
        error handling. Every protocol you meet answers one of them.
  \item Layers matter because they tell you what to test next. Work upwards from
        the physical layer; do not start at DNS.
  \item Switches forward frames using MAC addresses and flood when they do not
        know. Routers forward packets using IP addresses and drop when they do
        not know.
  \item One VLAN is one broadcast domain. Every switch port is its own collision
        domain.
  \item CCNA v2.0 has twenty-nine topics. Sixteen of them ask you to configure
        or to fix. Study accordingly.
\end{tight}
\end{chaptersummary}

\begin{reviewq}
  \item \textbf{(MCQ)} Which unit of data carries MAC addresses?
        (a)~segment (b)~packet (c)~frame (d)~datagram
  \item \textbf{(MCQ)} A device receives a frame for a destination MAC address
        it has never seen. What does it do?
        (a)~drops it (b)~floods it out all ports in the VLAN except the one it
        arrived on (c)~returns it to the sender (d)~routes it
  \item \textbf{(Short)} Give one reason the v2.0 blueprint dropped the
        requirement to describe the OSI model, and say what it expects instead.
  \item \textbf{(Short)} Explain the difference between a collision domain and a
        broadcast domain, and give the device that bounds each.
  \item \textbf{(Applied)} A colleague says ``we do not need to learn
        troubleshooting, it is only one section of the exam''. Using the verb
        counts in the blueprint, explain why that is wrong.
\end{reviewq}
